\documentclass[11pt]{scrartcl}
\usepackage[sexy]{evan} %BLBLBLBLBLBLBLBLBLBLBLB EVANCHEN
\usepackage[T1]{fontenc}
\usepackage[utf8]{inputenc}
\usepackage{graphicx}
\usepackage{amsmath}
\usepackage{tikz}
\usepackage{tikz-cd}
\usepackage{asymptote}
\usepackage{amsthm}
\usepackage{amssymb}
\usepackage{gensymb}
\usepackage{amsfonts}
\usepackage[english,italian]{babel}
\usepackage{quoting}
\quotingsetup{font=big}
\usepackage{booktabs}
\usepackage{caption}
\usepackage{lmodern}
\usepackage{faktor}
\usepackage{bbm} %oer fare gli insiemi indicatori
\usepackage{leftindex} %per gli ideali e gli indici a sinistra
\usepackage{hyperref}
\usepackage{epigraph} 
\usepackage{pgfplots}
\usepackage[export]{adjustbox}

%Pacchetti per i codici
\usepackage{listings}
\usepackage{xcolor}
\usepackage{algpseudocode}

%Analisi e probabilità
%----------------------------------
\newcommand{\dif}{\text{d}}
%%% the pure symbol, see https://tex.stackexchange.com/a/718194/4427
\newcommand{\cind}{\perp\mspace{-9mu}\perp}
%%% define \indep
\NewDocumentCommand{\indep}{e{_}}{%
	\mathrel{%
		\IfNoValueTF{#1}{% no subscript
			\cind
		}{% subscript
			\mathchoice{\mathop{\cind}\limits_{#1}}{\cind_{#1}}{\cind_{#1}}{\cind_{#1}}%
		}%
	}%
}
%-----------------------------------

%Algebra lineare e geometria
%----------------------------------
\newcommand{\tr}[1]{\text{tr}\!\left(#1\right)}
\renewcommand{\rank}[1]{\text{rank}\!\left(#1\right)}
\newcommand{\Immagine}[1]{\text{Im}\left(#1\right)}
\renewcommand{\ker}[1]{\text{Ker}\left(#1\right)}
\renewcommand{\AA}{\mathbb{A}}
\newcommand{\LL}{\mathbb{L}}
\newcommand{\SO}{\text{SO}}
\newcommand{\KK}{\mathbb{K}}
%----------------------------------

%Fisica
%----------------------------------
\newcommand{\UDiMisura}[1]{\,\text{#1}}
\newcommand{\ihat}{\hat{\imath}}
\newcommand{\jhat}{\hat{\jmath}}
\newcommand{\khat}{\hat{k}}
\newcommand{\posit}{\vec{r}}
\newcommand{\veloc}{\dot{\vec{r}}}
\newcommand{\accel}{\ddot{\vec{r}}}
%----------------------------------

%Informatica
%----------------------------------
\definecolor{codegreen}{rgb}{0,0.6,0}
\definecolor{codegray}{rgb}{0.5,0.5,0.5}
\definecolor{codepurple}{rgb}{0.58,0,0.82}
\definecolor{backcolour}{rgb}{0.95,0.95,0.92}

\lstdefinestyle{overleafwiki}{
	backgroundcolor=\color{backcolour},   
	commentstyle=\color{codegreen},
	keywordstyle=\color{magenta},
	numberstyle=\tiny\color{codegray},
	stringstyle=\color{codepurple},
	basicstyle=\ttfamily\footnotesize,
	breakatwhitespace=false,         
	breaklines=true,                 
	captionpos=b,                    
	keepspaces=true,                 
	numbers=left,                    
	numbersep=5pt,                  
	showspaces=false,                
	showstringspaces=false,
	showtabs=false,                  
	tabsize=2
}

\lstset{style=overleafwiki}
\newcommand{\bigo}[1]{\mathcal{O}\!\left(#1\right)}
\newcommand{\bigomega}[1]{\Omega\!\left(#1\right)}
\newcommand{\bigtheta}[1]{\Theta\!\left(#1\right)}
%----------------------------------

%Insiemistica e Algebra
%----------------------------------
\newcommand{\eset}{\varnothing}
\newcommand{\parti}[1]{\mathcal{P}\!\left(#1\right)}
\newcommand{\inj}{\hookrightarrow}
\newcommand{\surj}{\twoheadrightarrow}
\newcommand{\bij}{\hookrightarrow\mathrel{\mspace{-15mu}}\rightarrow}
\newcommand{\ordine}[2]{\text{ord}_{#1}\left(#2\right)}
\newcommand{\normal}{\mathrel{\unlhd}}
\newcommand{\MCD}{\text{MCD}}
\renewcommand{\mcm}{\text{mcm}}
%----------------------------------


\begin{document}
	\title{Appunti di Geometria 2}
	\subtitle{
		Orsola Tommasi (Modulo A) - UNIPD \\
		Nicola Mazzari (Modulo B) - UNIPD
	}
	\date{Settembre 2026 - Giugno 2027}
	
	
	\maketitle
	
	\tableofcontents
	
	\newpage
	
	\part{Modulo A}
	
	\section{Forme bilineari e quadratiche}
	
	Usiamo $\FF$ o $\KK$ per riferirci ad un campo generico.
	\begin{definition}[Applicazione multilineare]
		Dati $n\in\NN$ e spazi vettoriali $V_1,V_2,\dots,V_n,W$ su $\KK$, allora un'applicazione
		$$F:V_1\times V_2\times \dots \times V_n\to W$$
		è detta \textit{multilineare di grado $n$} o semplicemente \textit{$n$-lineare} se è lineare in ciascuna entrata, ovvero se, fissati $n-1$ argomenti, la funzione
		\begin{align*}
			F\left(v_1,v_2,\dots,v_{i-1},-,v_{i+1},\dots,v_n\right):V_i&\to  W \\
			v & \mapsto  F\left(v_1,v_2,\dots,v_{i-1},v,v_{i+1},\dots,v_n\right)
		\end{align*}
		è lineare per ogni $i$.
	\end{definition}
	\begin{example}[Prodotto scalare]
		Il primo esempio di tale funzione che conosciamo già è proprio il prodotto scalare standard:
		\begin{align*}
			\cdot:\RR^n\times \RR^n & \to \RR \\
			(v,w)&\mapsto v\cdot w = v^\intercal w = x_1y_1+\dots+x_ny_n.
		\end{align*}
	\end{example}
	
	\begin{example}[Determinante]
		Il determinante di una matrice $A$ può essere vista come una funzione multilineare delle colonne della matrice, che va da
		$$\det:\underset{\text{$n$ volte}}{\underbrace{\RR^n\times \RR^n \times \dots \times \RR^n}}\to\RR.$$
	\end{example}
	
	\begin{example}[Dualità canonica]
		La dualità canonica è naturalmente una funzione bilineare 
		$$\langle\cdot,\cdot\rangle:V^*\times V \to \KK.$$
	\end{example}
	
	\begin{example}[Funzione nulla]
		La funzione nulla su qualunque prodotto di spazi lineari è sempre multilineare.
	\end{example}
	
	Quello che tuttavia ci interessa studiare adesso sono le forme \textit{bilineari}:
	\begin{definition}[Forma bilineare]
		Una \textit{forma bilineare} è una funzione bilineare $B:V\times V\to \KK$, dove $V$ e $W$ sono spazi $\KK$-lineari.
	\end{definition}

	\begin{definition}[Simmetria e antisimmetria]
		Una forma bilineare $B(v,w)$ si dice \textit{simmetrica} se, per ogni $v,w\in V$
		$$B(v,w)=B(w,v)$$
		e \textit{antisimmetrica} se per ogni $v,w\in V$:
		$$B(v,w)=-B(w,v)$$
		infine si dice \textit{alternante} se per ogni $v\in V$ si ha
		$$B(v,v)=0$$
		com'è noto dal caso del determinante.
	\end{definition}
	La forma bilineare nasce come generalizzazione naturale della nozione di prodotto scalare, e quindi è chiaro che il prodotto scalare standard sia una forma bilineare simmetrica. Un esempio più elaborato di forma bilineare è
	\begin{example}[Prodotto in metriche non euclidee]
		Sia $A\in M_n(\KK)$ una matrice fissata con coefficienti $a_{ij}$, allora la funzione
		\begin{align*}
			g_A:\KK^n\times \KK^n & \to \KK \\
			(u,v) & \mapsto u \cdot Av = u^\intercal A v=\sum_{1\le i,j \le n}a_{ij}x_iy_j
		\end{align*}
		è sempre una forma bilineare, che soddisfa 
		$$a_{ij}=g_{A}\left(e_i,e_j\right).$$
	\end{example}
	\begin{example}[Prodotto in un campo]
		Il prodotto in qualunque campo $\KK$ è una forma bilineare simmetrica e non alternante. Se inoltre il campo ha caratteristica $2$, questa forma è anche antisimmetrica, in quanto $1=-1$.
	\end{example}
	\begin{example}[Determinante]
		La funzione
		\begin{align*}
			B:\KK^2\times \KK^2 & \to \KK \\
			\begin{pmatrix}
				x_1 \\
				y_2 
			\end{pmatrix},
			\begin{pmatrix}
				x_2 \\
				y_2
			\end{pmatrix} & \mapsto \det
			\begin{pmatrix}
				x_1 & x_2 \\
				y_1 & y_2
			\end{pmatrix}
		\end{align*}
		è bilineare alterna, ma cosa succede se identifichiamo $\KK$ con $\RR$ e $\KK^2$ con $\CC$?
	\end{example}
	
	\subsection{Prime proprietà e classificazione grossolana}
	
	Vogliamo ora capire almeno un minimo come si classifichino le forme bilineari,
	il primo risultato che dimostriamo è
	\begin{proposition}[Relazione tra simmetria, antisimmetria e alternanza]
		Abbiamo le seguenti proposizioni:
		\begin{itemize}
			\item Su ogni campo $\KK$, ogni forma alternante è antisimmetrica.
			\item Sui campi a caratteristica diversa da $2$, una forma bilineare è alternante se e solo se è antisimmetrica.
			\item Sui campi a caratteristica $2$, una forma bilineare è simmetrica se e solo se è antisimmetrica.
		\end{itemize}
	\end{proposition}
	\begin{proof}
		Dimostriamo una cosa alla volta:
		\begin{itemize}
			\item Per alternanza abbiamo
			$$B(v+w,v+w)=0$$
			d'altra parte espandendo per bilinearità e cancellando $B(v,v)$ e $B(w,w)$ per alternanza:
			$$B(v,w)+B(w,v)=0$$
			da cui la tesi.
			\item Per antisimmetria abbiamo
			$$B(v,v)=-B(v,v)$$
			cioè
			$$2B(v,v)=0$$
			da cui, dato che la caratteristica è diversa da $2$, abbiamo la tesi.
			\item Semplicemente segue dal fatto che $1=-1$.
		\end{itemize}
		Abbiamo dimostrato tutti i punti, quindi abbiamo finito.
	\end{proof}
	
	\begin{theorem}[di decomposizione]
		Se $\text{car}(\KK)\ne 2$ allora ogni forma bilineare $B$ definita su $\KK$ si scrive in modo unico come somma di una forma simmetrica e di una alternante.
	\end{theorem}
	\begin{proof}
		Per l'esistenza, possiamo pensare euristicamente di star cercando due forme $B_S$ e $B_A$ tali che
		$$
		\begin{cases}
			B(v,w)=B_S(v,w)+B_A(v,w) \\
			B(w,v)=B_S(v,w)-B_A(v,w)
		\end{cases}
		$$
		da cui possiamo ricavare direttamente delle formule per le due forme cercate
		\begin{gather*}
			B_S(v,w)=\frac{B(v,w)+B(w,v)}{2} \\
			B_A(v,w)=\frac{B(v,w)-B(w,v)}{2}
		\end{gather*}
		e visto che abbiamo trovato una formula esplicita, questo ci dice anche che la decomposizione è unica.
	\end{proof}
	
	Possiamo dunque definire i seguenti spazi di funzioni:
	\begin{definition}[Spazi bilineari]
		Definiamo i seguenti spazi
		\begin{gather*}
			\text{Bil}_\KK(V):=\left\{\text{forme bilineari su $V$}\right\} \\
			\text{Alt}_\KK(V):=\left\{B\in \text{Bil}_\KK(V):\text{$B$ è alternante}\right\} \\
			\text{Antisimm}_\KK(V):=\left\{B\in \text{Bil}_\KK(V):\text{$B$ è antisimmetrica}\right\} \\
			\text{Simm}_\KK(V):=\left\{B\in \text{Bil}_\KK(V):\text{$B$ è simmetrica}\right\}
		\end{gather*}
		che si dimostra agilmente essere tutti spazi vettoriali.
	\end{definition}
	Possiamo quindi riformulare quanto visto finora in termini di sottospazi vettoriali
	\begin{corollary}[di Decomposizione]
		Per campi $\KK$ con caratteristica diversa da $2$ abbiamo
		$$\text{Simm}_\KK(V)\oplus \text{Alt}_\KK(V)=\text{Bil}_\KK(V)$$
		mentre se la caratteristica è esattamente $2$:
		$$\text{Alt}_\KK(V)\subset \text{Antisimm}_\KK(V)=\text{Simm}_\KK(V)$$
	\end{corollary}
	Tiriamo in ballo le basi con il seguente risultato:
	\begin{theorem}[Informazione sufficiente per determinare forme simmetriche]
		Se $\KK$ è un campo a caratteristica diversa da $2$ allora una forma bilineare simmetrica $B$ su $V$ è determinata in modo univoco dal suo valore per le coppie $(v,v)$ con $v\in V$.
	\end{theorem}
	\begin{proof}
		Semplicemente dall'identità, valida solo per le forme simmetriche:
		$$B(v+w,v+w)=B(v,v)+2B(v,w)+B(w,w)$$
		otteniamo che
		$$B(v,w)=\frac{B(v+w,v+w)-B(v,v)-B(w,w)}{2}$$
		per cui appunto ci basta conoscere solo termini della forma $B(u,u)$, con $u\in V$. Con questa formula possiamo calcolare qualunque valore di $B$, e quindi abbiamo finito.
	\end{proof}
	Purtroppo questo teorema ci richiede di conoscere \textbf{tutti} i prodotti $B(v,v)$, e non solo quelli su una base, come mostra il seguente esempio:
	\begin{example}[Una forma non caratterizzata dalla forma quadratica sulla base]
		La forma
		$$B:\KK^2\times \KK^2\to \KK$$
		tale che
		$$B\left(
		\begin{pmatrix}
			x_1 \\
			y_1
		\end{pmatrix},
		\begin{pmatrix}
			x_2 \\
			y_2
		\end{pmatrix}
		\right)=x_1y_2+x_2y_1$$
		per cui abbiamo che $B(e_1,e_1)=B(e_2,e_2)=0$, che chiaramente non basta per identificare tutti i prodotti $B(v,w)$.
	\end{example}
	Facciamo qualche altro esempio di forma bilineare:
	\begin{example}[Prodotto scalare negativo]
		Definiamo il prodotto scalare
		$$B\left(
		\begin{pmatrix}
			x_1 \\
			y_1
		\end{pmatrix},
		\begin{pmatrix}
			x_2 \\
			y_2
		\end{pmatrix}
		\right)=x_1x_2-y_1y_2$$
	\end{example}
	\begin{example}[Prodotto misto]
		Sia $u\in \KK^3$, definiamo allora
		\begin{align*}
			B_u:\KK^3\times \KK^3 & \to \KK^3 \\
			(v,w) & \mapsto \det\left(
			\begin{array}{c|c|c}
				v & w & u
			\end{array}
			\right)=\left(v\times w\right)\cdot u
		\end{align*}
		Questo è il cosiddetto prodotto misto, ed è sia antisimmetrica che alternante.
	\end{example}
	In effetti possiamo dimostrare il seguente risultato più forte:
	\begin{proposition}[Rappresentazione delle forme alternanti su $\KK^3$]
		Tutte le forme bilineari alternati su $\KK^3$ sono della forma
		$$B(v,w)=\left(v\times w\right)\cdot u$$
		per qualche $u\in V$.
	\end{proposition}
	
	\begin{example}[Duale]
		Un altro esempio di forma bilineare è il seguente:
		\begin{align*}
			B:\left(V\oplus V^*\right)\times \left(V\oplus V^*\right) & \to \KK \\
			\left(\left(\varphi,v\right),\left(\psi,w\right)\right) & \mapsto \psi(v)-\varphi(w)
		\end{align*}
		e questa è chiaramente alternante. La notazione $V\oplus V^*$ è semplicemente un modo alternativo per scrivere $V\times V^*$, solo che al posto di considerare le coppie $(v,w)$ consideriamo le "somme formali" $v+w$ tra elementi di $V$ e $W$.
	\end{example}
	\begin{example}[Funzionali bilineari in analisi]
		Se consideriamo lo spazio vettoriale $V=C^0\left([0,1]\right)$ delle funzioni continue su $[0,1]$, allora possiamo definire una forma bilineare come
		$$B(f,g)=\int_0^1 f(x)g(x)\dif x$$
		per ogni $f,g\in V$.
	\end{example}
	\begin{example}[Forma Hermitiana]
		Prendiamo lo spazio $V=\CC^n$ e definiamo la forma
		$$H
		\left(\begin{pmatrix}
			z_1 \\
			\vdots \\
			z_n 
		\end{pmatrix}
		,
		\begin{pmatrix}
			w_1 \\
			\vdots \\
			w_n 
		\end{pmatrix}\right)=\sum_{i=1}^{n}z_i\overline{w}_i
		$$ 
		e qui il campo su cui definiamo $V$ è importante! Su $\RR$ infatti $H$ risulta bilineare, ma su $\CC$ abbiamo che
		$$H(v,\lambda w)=\overline{\lambda}H(v,w)$$
		e questa è proprio la proprietà che vedremo avrà una forma hermitiana, in particolare questa è la cosiddetta forma hermitiana standard.
	\end{example}
	
	\subsubsection{La traccia}
	
	Introduciamo ora una funzione estremamente importante
	
	\begin{definition}[Traccia]
		Data una matrice $A\in M_n(\KK)$ di coefficienti $a_{ij}$ definiamo la sua \textit{traccia} come
		$$\tr A:=\sum_{i=1}^{n}a_{ii}$$
		cioè la somma degli elementi sulla diagonale.
	\end{definition}
	
	\begin{example}[Traccia]
		Possiamo definire una forma bilineare simmetrica usando la traccia:
		\begin{align*}
			B:M_n(\KK)\times M_n(\KK) & \to \KK \\
			(A,A') & \mapsto \tr{AA'}.
		\end{align*}
	\end{example}
	
	La simmetria della forma descritta sopra deriva dal seguente risultato:
	\begin{lemma}[Ciclicità della traccia]
		Per ogni $A,B\in M_n(\KK)$ abbiamo che
		$$\tr{AB}=\tr{BA}.$$
	\end{lemma}
	\begin{proof}
		Calcolo diretto.
	\end{proof}
	
	INSERISCI QUALCOSA SUI COMMUTATORI/QUALCOSA DALL'AXLER
	
	\subsection{Ortogonalità}
	
	Per il prodotto scalare una grande parte della teoria consisteva nella nozione di ortogonalità, riprendiamola qui dunque come:
	\begin{definition}[Ortogonalità]
		Se $B$ è una forma bilineare su $V$ allora diciamo che $v\in V$ è \textit{ortogonale} a $w$ se $B(v,w)=0$, e in questo caso scriviamo
		$$v\perp_B w$$
		oppure quando la forma di cui stiamo parlando è chiara, $v\perp w$.
	\end{definition}
	La prima cosa strana che notiamo è che per forme arbitrarie l'ortogonalità può non essere simmetrica, ovvero
	\begin{example}[Una nozione di ortogonalità non simmetrica]
		Se consideriamo la forma su $\RR^2$ definita come
		$$B(\mathbf{x},\mathbf{y})=x_1y_1+x_2y_1+x_2y_2$$
		abbiamo che $e_1\perp_B e_2$ e $e_2\perp_B e_1$.
	\end{example}
	Per caratterizzare quando la perpendicolarità sia simmetrica allora abbiamo bisogno di questo lemma:
	\begin{lemma}[Condizione necessaria per la simmetria della perpendicolarità]
		Se $\perp_B$ è una relazione simmetrica allora per ogni $u,v,w\in V$ si ha
		$$B(w,u)B(u,v)=B(v,u)(u,w).$$
	\end{lemma}
	\begin{proof}
		Il sottospazio generato da $v,w$ contiene il vettore 
		$$z=B(w,u)v-B(v,u)w$$
		e a questo punto notiamo che $B(z,u)=0$, quindi visto che abbiamo assunto che l'ortogonalità è simmetrica abbiamo
		$$B(u,z)=0$$
		che espansa è proprio la tesi.
	\end{proof}
	\begin{theorem}[Caratterizzazione della simmetria dell'ortogonalità]
		Se $B$ è una forma bilineare su $V$ allora $\perp_B$ è una relazione simmetrica se e solo se $B$ è simmetrica o alternante.
	\end{theorem}
	\begin{proof}
		Un verso è chiaro. Per dimostrare l'altro verso 
		usiamo il lemma sopra con $u=w$, e abbiamo
		$$B(u,u)B(u,v)=B(v,u)B(u,u)$$
		cioè
		$$\left[B(u,v)-B(v,u)\right]B(u,u)=0$$
		che vuol dire che per ogni coppia di vettori $B$ si comporta o come se fosse alternante o come se fosse simmetrica.
	\end{proof}
	
	E quindi possiamo definire:
	\begin{definition}[Sottospazio ortogonale]
		Dato un $S\subseteq V$ possiamo definire il suo \textit{sottospazio ortogonale} $S^{\perp_B}$ come lo spazio dei vettori ortogonali a ciascun elemento di $S$. Si dimostra facilmente che questo è un sottospazio, e se per un $W$ spazio vettoriale abbiamo anche che
		$$W\oplus W^{\perp_B}=V$$
		questo sottospazio si dice \textit{complemento ortogonale} di $W$. 
	\end{definition}
	\begin{definition}[Somma diretta e restrizioni]
		Sia $B$ una forma bilineare su $V$. Allora per ogni sottospazio $W\le V$ possiamo definire una forma bilineare ristretta
		$$B_{\mid W}:W\times W \to \KK$$
		che definiremo \textit{forma bilineare indotta}. Allo stesso modo, dati due sottospazi di $V_1$ e $V_2$ tali che
		$$V=V_1\oplus V_2$$
		e due forme bilineari $B_1$ e $B_2$ definite rispettivamente su $V_1$ e $V_2$ possiamo definire una forma \textit{somma diretta} delle due, come:
		\begin{align*}
			B_1\oplus B_2 : V\times V & \to \KK \\
			(v_1+v_2,v_1'+v_2') & \mapsto B_1(v_1,v_1')+B_2(v_2,v_2')
		\end{align*}
		Questa costruzione, lo spazio $(V_1\oplus V_2,B_1\oplus B_2)$, viene detta la \textit{somma diretta ortogonale} di $V_1$ e $V_2$ e si indica come $V_1\perp V_2:= V$.
	\end{definition}
	Un esempio estremamente naturale di ciò è il seguente:
	\begin{example}[Prodotto scalare]
		Se consideriamo $V_1=V_2=\left(\RR,\cdot\right)$ possiamo effettuare la loro somma diretta ortogonale $\RR\perp\RR$, ottenendo esattamente uno spazio $(\RR^2,B)$, dove $B$ è definita come
		$$B(\mathbf{x},\mathbf{y})=x_1y_1+x_2y_2$$
		che è proprio il prodotto scalare standard! 
	\end{example}
	\begin{proposition}[Caratterizzazione della somma diretta ortogonale]
		Se $V=V_1\oplus V_2$ e $B_1,B_2$ sono forme bilineari rispettivamente su $V_1$ e $V_2$ allora $B_1\oplus B_2$ è l'unica forma bilineare su $V$ tale che:
		\begin{itemize}
			\item Le restrizioni sono le forme di partenza, ovvero $B_{\mid V_1}=B_1$ e $B_{\mid V_2}=B_2$. 
			\item Per ogni $v_1\in V_1$ e $v_2\in V_2$ si ha $B(v_1,v_2)=0$.
		\end{itemize}
	\end{proposition}
	\begin{proof}
		contenuto...
	\end{proof}
	
	\newpage
	
	\part{Modulo B}
	
\end{document}
